\section{Описание предметной области и структуры данных}

Для обучения и тестирования интеллектуального модуля была выбрана демонстрационная база данных «Авиаперевозки» (схема \texttt{bookings}), предоставленная компанией Postgres Pro. Данная база является репрезентативной для задач класса Text-to-SQL, поскольку она моделирует реальные бизнес-процессы авиакомпании и сочетает строго структурированные реляционные данные с полуструктурированными элементами в формате \texttt{JSONB}. Это позволяет одновременно оценивать как корректность синтаксической генерации SQL-запросов, так и способность системы интерпретировать вложенные структуры данных.

\subsection{Обоснование выбора и источник данных}

Выбор данной базы данных обусловлен ее архитектурной сложностью, объемом и особенностями хранения данных. Схема включает восемь взаимосвязанных таблиц, что создает реалистичную среду для проверки способности интеллектуального агента строить корректные соединения между сущностями. Используемый набор данных \texttt{demo-medium}, содержащий информацию за один год, формирует достаточную нагрузку для анализа производительности и устойчивости генерируемых запросов. Дополнительным фактором является наличие мультиязычных атрибутов, реализованных через \texttt{JSONB}, что требует от модели понимания структуры вложенных ключей и контекста пользовательского запроса.

\subsection{Логическая структура базы данных}

Логическая модель базы данных организована в виде нескольких взаимосвязанных функциональных блоков, каждый из которых отражает отдельный аспект предметной области. Основные сущности и их назначение представлены в таблице~\ref{tab:erd_blocks}.

\begin{table}[h!]
\centering
\caption{Функциональные блоки логической модели}
\label{tab:erd_blocks}
\small
\begin{tabularx}{\textwidth}{|>{\raggedright\arraybackslash}p{0.20\textwidth}|>{\raggedright\arraybackslash}p{0.32\textwidth}|>{\raggedright\arraybackslash}X|}
\hline
\textbf{Блок} & \textbf{Таблицы} & \textbf{Назначение} \\
\hline
Бронирование & \texttt{bookings}, \texttt{tickets} & Хранение информации о бронированиях и пассажирах \\
\hline
Рейсы & \texttt{flights}, \texttt{routes} & Описание расписания, статусов и маршрутов \\
\hline
Обслуживание & \texttt{ticket\_flights}, \texttt{boarding\_passes}, \texttt{seats} & Связь билетов с рейсами и управление посадкой \\
\hline
Справочники & \texttt{aircrafts\_data}, \texttt{airports\_data} & Метаданные по самолетам и аэропортам \\
\hline
\end{tabularx}
\end{table}

Такая декомпозиция упрощает интерпретацию схемы и позволяет интеллектуальному модулю выделять контекст запроса, определяя, к какому функциональному блоку относится пользовательский вопрос.

\subsection{Специфика хранения метаинформации}

Особенностью базы данных является использование типа \texttt{JSONB} для хранения мультиязычных атрибутов, что особенно важно для справочных таблиц. Например, в таблице \texttt{aircrafts\_data} поле \texttt{model} содержит названия воздушных судов сразу на нескольких языках. Извлечение значений требует явного указания ключа, соответствующего языку запроса. Структура хранения иллюстрируется в таблице~\ref{tab:json_example}.

\begin{table}[h!]
\centering
\caption{Пример структуры JSONB-поля}
\label{tab:json_example}
\begin{tabular}{|l|l|}
\hline
\textbf{Ключ} & \textbf{Значение} \\
\hline
\texttt{model -> 'ru'} & Сессна 208 Караван \\
\hline
\texttt{model -> 'en'} & Cessna 208 Caravan \\
\hline
\end{tabular}
\end{table}

Данная особенность накладывает дополнительные требования на этап генерации запросов, поскольку система должна корректно определять язык пользовательского ввода и выбирать соответствующий ключ внутри JSON-структуры.

\subsection{Фрагменты данных}

Для тестирования интеллектуального модуля использовались реальные фрагменты данных. Пример содержимого таблицы \texttt{aircrafts\_data} приведен в таблице~\ref{tab:aircrafts}, где представлены коды самолетов, их модели и дальность полета.

\begin{table}[h!]
\centering
\caption{Фрагмент справочника воздушных судов}
\label{tab:aircrafts}
\small
\begin{tabularx}{\textwidth}{|>{\raggedright\arraybackslash}p{0.20\textwidth}|>{\raggedright\arraybackslash}X|>{\centering\arraybackslash}p{0.12\textwidth}|}
\hline
\textbf{aircraft\_code} & \textbf{model (JSONB)} & \textbf{range} \\
\hline
773 & \{"en": "Boeing 777-300", "ru": "Боинг 777-300"\} & 11100 \\
\hline
SU9 & \{"en": "Sukhoi Superjet-100", "ru": "Сухой Суперджет-100"\} & 3000 \\
\hline
\end{tabularx}
\end{table}

Аналогично, таблица \texttt{flights} содержит информацию о конкретных рейсах, включая расписание и статус выполнения, что показано в таблице~\ref{tab:flights}.

\begin{table}[h!]
\centering
\caption{Фрагмент данных о рейсах}
\label{tab:flights}
\small
\begin{tabularx}{\textwidth}{|>{\centering\arraybackslash}p{0.12\textwidth}|>{\centering\arraybackslash}p{0.12\textwidth}|>{\raggedright\arraybackslash}X|>{\centering\arraybackslash}p{0.15\textwidth}|>{\centering\arraybackslash}p{0.16\textwidth}|}
\hline
\shortstack{\textbf{flight}\\\textbf{id}} & \shortstack{\textbf{flight}\\\textbf{no}} & \shortstack{\textbf{scheduled}\\\textbf{departure}} & \shortstack{\textbf{departure}\\\textbf{airport}} & \textbf{status} \\
\hline
1391 & PG0485 & 2026-05-01 06:30:00 & DME & Scheduled \\
\hline
2540 & PG0021 & 2026-05-01 12:00:00 & VKO & Arrived \\
\hline
\end{tabularx}
\end{table}

\subsection{Подготовка данных для RAG-конвейера}

Для интеграции базы данных в RAG-конвейер была выполнена предварительная текстовая интерпретация структуры каждой таблицы. Вместо прямой передачи схемы в модель используется подход семантического индексирования, при котором каждая таблица представляется в виде компактного текстового описания. Пример такого представления приведен в таблице~\ref{tab:rag_doc}.

\begin{table}[h!]
\centering
\caption{Пример документа для векторного хранилища}
\label{tab:rag_doc}
\begin{tabular}{|p{3cm}|p{10cm}|}
\hline
\textbf{Атрибут} & \textbf{Содержание} \\
\hline
Таблица & \texttt{flights} \\
\hline
Описание & Содержит информацию о рейсах и их статусах \\
\hline
Ключевые поля & \texttt{flight\_id}, \texttt{status}, \texttt{departure\_airport} \\
\hline
Связи & Соединение с \texttt{airports\_data} по коду аэропорта \\
\hline
\end{tabular}
\end{table}

Подобный подход позволяет интеллектуальному агенту динамически извлекать релевантную информацию о структуре базы данных, минимизируя объем передаваемого контекста и повышая точность генерации SQL-запросов. Это особенно важно в условиях ограничений на размер входных данных для языковых моделей.
